% Informe in situ - Sesión 3: Oscilador sinusoidal en puente de Wien
% Instrumentación Electrónica - Grado en Física - Universidad de Salamanca
\documentclass[11pt,a4paper]{article}

\usepackage[utf8]{inputenc}
\usepackage[T1]{fontenc}
\usepackage[spanish,es-noshorthands]{babel}
\usepackage{mathpazo}
\usepackage{microtype}
\usepackage[a4paper,top=1.2cm,bottom=1.2cm,left=2cm,right=2cm]{geometry}
\usepackage{xcolor}
\usepackage{booktabs}
\usepackage{array}
\usepackage{etoolbox}
\usepackage{siunitx}
\usepackage[siunitx,europeanresistors,americaninductors]{circuitikz}
\usepackage{titlesec}
\usepackage{enumitem}

\sisetup{output-decimal-marker={,}, group-minimum-digits=4}

\definecolor{azul}{HTML}{1F3A5F}
\definecolor{gris}{gray}{0.45}
\definecolor{fondo}{HTML}{EEF2F7}

\titleformat{\section}{\large\bfseries\color{azul}}{\thesection.}{0.5em}{}
\titlespacing*{\section}{0pt}{0.7em}{0.25em}
\setlength{\parindent}{0pt}
\setlength{\parskip}{0.35em}
\pagestyle{plain}

% ---------------------------------------------------------------------
% Datos medidos. Dejar vacíos para rellenar a mano en el laboratorio;
% si se escribe un valor, se imprime en lugar del hueco.
% ---------------------------------------------------------------------
\newcommand{\Rmuno}{}   % R1 medida (kΩ)
\newcommand{\Rmdos}{}   % R2 medida (kΩ)
\newcommand{\Rmtres}{}  % R3 medida (kΩ)
\newcommand{\Rmcuatro}{}% R4 medida (kΩ)
\newcommand{\Cmuno}{}   % C1 medido (nF)
\newcommand{\Cmdos}{}   % C2 medido (nF)
\newcommand{\Rpot}{}    % R4 (potenciómetro) tras el ajuste (kΩ)

% Hueco para rellenar: \hueco[ancho]{valor}
\newcommand{\hueco}[2][1.8cm]{%
  \expandafter\ifblank\expandafter{#2}{\rule[-0.25ex]{#1}{0.4pt}}{\makebox[#1]{#2}}}
% Líneas en blanco para respuestas: \lineas{n}
\newcommand{\lineas}[1]{%
  \par\vspace{-0.4em}%
  \foreach \i in {1,...,#1}{\vspace{1.05em}\noindent\rule{\linewidth}{0.3pt}\par}}

\begin{document}

% ------------------------- Cabecera ----------------------------------
{\large\bfseries\color{azul} Sesión 3: Oscilador sinusoidal en puente de Wien}\hfill
{\small\color{gris} Instrumentación Electrónica $\cdot$ Grado en Física}\\[0.2em]
{\small Miguel Martín Cortés (70974284X) $\cdot$ Álvaro Montero Mateos (71096223A) $\cdot$
Daniel Martín Guerrero (71063727Y)}\\[0.1em]
{\small Grupo C.2 \hfill Fecha: \hueco[2.2cm]{} de noviembre de 2026}\\[-0.5em]
{\color{azul}\rule{\linewidth}{0.5pt}}

% ------------------------- Introducción ------------------------------
\section{Introducción}

\noindent
\begin{minipage}[c]{0.58\textwidth}
Se monta un oscilador sinusoidal en puente de Wien con el amplificador
operacional \textmu A741 (figura~1). La red $RC$ de realimentación positiva
selecciona la frecuencia de oscilación, y la red $R_3$--$R_4$ de
realimentación negativa fija la ganancia $K = 1 + R_4/R_3$. Con
$R_1 = R_2 = R$ y $C_1 = C_2 = C$, las condiciones de oscilación son
\[
  R_4 = 2R_3 \;\;(K = 3), \qquad f_0 = \frac{1}{2\pi RC}.
\]
Se comprueba la forma de onda de salida, se ajusta $R_4$ con un
potenciómetro hasta obtener una sinusoide limpia y se mide $f_0$ para
varios pares de condensadores.
\end{minipage}\hfill
\begin{minipage}[c]{0.38\textwidth}
\centering
\begin{circuitikz}[scale=0.58, transform shape, font=\large]
  \draw (0,0) node[op amp, noinv input up] (oa) {};
  \draw (oa.out) -- ++(0.8,0) coordinate (vo) node[right] {$v_o$};
  % Realimentación positiva: R1--C1 en serie, R2 || C2 a masa
  \draw (oa.+) -- ++(-3.0,0) coordinate (vp);
  \draw (vp) -- (vp |- 0,2.2) coordinate (t1)
        to[C, l=$C_1$] ($(t1)!0.5!(vo |- 0,2.2)$)
        to[R, l=$R_1$] (vo |- 0,2.2) -- (vo);
  \draw (vp) to[short, *-] ++(0,-2.4) coordinate (a);
  \draw (a) -- ++(-0.8,0) to[R, l_=$R_2$] ++(0,-2.0) -- ++(0.8,0) coordinate (b);
  \draw (a) -- ++(0.8,0) to[C, l=$C_2$] ++(0,-2.0) -- ++(-0.8,0);
  \draw (b) node[ground] {};
  % Realimentación negativa: R4 a la salida, R3 a masa
  \draw (oa.-) -- ++(-0.4,0) coordinate (m) -- (m |- 0,-2.0) coordinate (vm);
  \draw (vm) to[R, l_=$R_4$, *-] (vo |- vm) -- (vo);
  \draw (vo) node[circ] {};
  \draw (vm) to[R, l=$R_3$] ++(0,-2.4) node[ground] {};
\end{circuitikz}\\[0.2em]
{\small\color{gris} Figura 1. Oscilador en puente de Wien\\(sin la polarización $\pm V_{CC}$).}
\end{minipage}

% ------------------------- Componentes -------------------------------
\section{Valores de los componentes}

\begin{center}
\renewcommand{\arraystretch}{1.3}
\begin{tabular}{l*{6}{>{\centering\arraybackslash}p{1.45cm}}}
  \toprule
  & $R_1$ [\si{\kilo\ohm}] & $R_2$ [\si{\kilo\ohm}] & $R_3$ [\si{\kilo\ohm}]
  & $R_4$ [\si{\kilo\ohm}] & $C_1$ [\si{\nano\farad}] & $C_2$ [\si{\nano\farad}] \\
  \midrule
  Nominal & 10 & 10 & 10 & 20 & 10 & 10 \\
  Medido (multímetro) & \hueco[1.2cm]{\Rmuno} & \hueco[1.2cm]{\Rmdos} & \hueco[1.2cm]{\Rmtres}
          & \hueco[1.2cm]{\Rmcuatro} & \hueco[1.2cm]{\Cmuno} & \hueco[1.2cm]{\Cmdos} \\
  \bottomrule
\end{tabular}
\end{center}

% ------------------------- Forma de onda -----------------------------
\section{Forma de onda a la salida}

\textbf{¿Qué onda se mide en la salida con el osciloscopio?}
\lineas{2}

\vspace{0em}
\textbf{Ajuste con potenciómetro.} Sustituida $R_4$ por un potenciómetro
(ajustado inicialmente a $\approx\SI{20}{\kilo\ohm}$) y retocado hasta que la
salida se parece lo más posible a una sinusoide:
\[
  R_4 = \hueco[2.2cm]{\Rpot}\ \si{\kilo\ohm},
  \qquad
  R_4/R_3 = \hueco[2.2cm]{}.
\]

% ------------------------- Frecuencia --------------------------------
\section{Frecuencia de oscilación}

Con el potenciómetro reajustado para mantener la salida sinusoidal, se
sustituye el par de condensadores y se miden automáticamente con el
osciloscopio el periodo $T_0$ y la frecuencia $f_0$.

\begin{center}
\renewcommand{\arraystretch}{1.3}
\begin{tabular}{*{6}{>{\centering\arraybackslash}p{2.05cm}}}
  \toprule
  $C$ [\si{\pico\farad}] & $R$ [\si{\kilo\ohm}] & $f_0$ [\si{\kilo\hertz}] calculada
  & $T_0$ [\si{\milli\second}] medido & $f_0$ [\si{\kilo\hertz}] medida
  & Desviación [\%] \\
  \midrule
  & 10 & & & & \\
  & 10 & & & & \\
  & 10 & & & & \\
  \bottomrule
\end{tabular}
\end{center}

\textbf{¿Se ajusta $f_0$ al valor obtenido con la ecuación $f_0 = 1/(2\pi RC)$?}
\lineas{3}

\end{document}
